\section*{Open Questions}
\textbf{Q1.} What is the concrete size overhead of Brakerski-Yuen's Clifford-only QGC as
a function of circuit depth $d$ and wire count $w$? The paper gives asymptotic bounds
(double-exponential in $d$ information-theoretically; polynomial with QPRGs) but not a
tight constant; a small-instance simulator that emits the explicit garbling and counts
EPR pairs $+$ classical bits would give practitioners a real-world dial. \emph{Basis:}
our toy replication needed only a handful of EPR pairs per gate, but the paper's proof
sketch suggests the constant hidden inside ``polynomial'' could be large.

\textbf{Q2.} How does the Clifford-slice Pauli-frame propagation degrade under realistic
noise? Our numerical checks used ideal density matrices; a natural follow-up is to
inject depolarizing noise on the EPR pairs and quantify the decoded-state fidelity drop
as a function of physical error rate --- i.e., is QGC ``fault-tolerant-friendly'' in the
same sense as Steane-code teleportation? \emph{Basis:} the paper is noise-free by
construction; but any deployed QGC will inherit teleportation's noise sensitivity.

\textbf{Q3.} Can the double-exponential-in-depth blow-up of the information-theoretic
variant be replaced by a doubly-exponential-in-\emph{$T$-count} blow-up? The paper
identifies non-Clifford gates as the source of the blow-up but does not sharpen the
bound to explicitly track $T$-count separately from depth. \emph{Basis:} our Clifford
slice had zero blow-up; the entire cost seems to sit on the $T$-gate magic-state gadget.

\textbf{Q4.} What is the minimal quantum-secure PRG family sufficient for the
computational variant of QGC, and does a lattice-based candidate (e.g.\ SIS-based
$G_\lambda$) yield a concrete instantiation with practical parameters (say, security
$\lambda=128$, gate count $<10^4$)? \emph{Basis:} the paper cites ``QPRGs'' abstractly;
a concrete choice would let one benchmark end-to-end QGC latency for the first time.

\textbf{Q5.} The paper's $\Sigma$-protocol for QMA relies on QGC as a black box; does
the constant-depth encoder property translate into a constant-round QMA protocol whose
prover runs in $\mathsf{QNC}^0$? If so, this would be a strictly stronger complexity
statement than the paper's stated round complexity. \emph{Basis:} we did not exercise
the protocol layer, but the encoder's $\mathsf{QNC}^0_f$ property is the natural
resource-count to inherit into the higher-level protocol.
